\section{Preparation for tame symbols}
\label{section-preparation-tame-symbol}

\noindent
In this section we put some lemma that will help us define the
tame symbol in the next section.

\begin{lemma}
\label{lemma-glue-at-max}
Let $A$ be a Noetherian ring. Let $\mathfrak m_1, \ldots, \mathfrak m_r$
be pairwise distinct maximal ideals of $A$. For $i = 1, \ldots, r$ let
$\varphi_i : A_{\mathfrak m_i} \to B_i$ be a ring map whose
kernel and cokernel are annihilated by a power
of $\mathfrak m_i$. Then there exists a ring map $\varphi : A \to B$ such
that
\begin{enumerate}
\item the localization of $\varphi$ at $\mathfrak m_i$ is
isomorphic to $\varphi_i$, and
\item $\Ker(\varphi)$ and $\Coker(\varphi)$ are annihilated
by a power of $\mathfrak m_1 \cap \ldots \cap \mathfrak m_r$.
\end{enumerate}
Moreover, if each $\varphi_i$ is finite, injective, or
surjective then so is $\varphi$.
\end{lemma}

\begin{proof}
Set $I = \mathfrak m_1 \cap \ldots \cap \mathfrak m_r$. Set
$A_i = A_{\mathfrak m_i}$ and $A' = \prod A_i$.
Then $IA' = \prod \mathfrak m_i A_i$ and $A \to A'$
is a flat ring map such that $A/I \cong A'/IA'$.
Thus we may use More on Algebra, Lemma
\ref{more-algebra-lemma-application-formal-glueing}
to see that there exists an $A$-module map $\varphi : A \to B$
with $\varphi_i$ isomorphic to the localization of $\varphi$
at $\mathfrak m_i$. Then we can use the discussion in
More on Algebra, Remark \ref{more-algebra-remark-formal-glueing-algebras}
to endow $B$ with an $A$-algebra structure
matching the given $A$-algebra structure on $B_i$.
The final statement of the lemma follows easily from
the fact that $\Ker(\varphi)_{\mathfrak m_i} \cong \Ker(\varphi_i)$
and $\Coker(\varphi)_{\mathfrak m_i} \cong \Coker(\varphi_i)$.
\end{proof}

\noindent
The following lemma is very similar to
Algebra, Lemma \ref{algebra-lemma-nonregular-dimension-one}.

\begin{lemma}
\label{lemma-Noetherian-domain-dim-1-two-elements}
Let $(R, \mathfrak m)$ be a Noetherian local ring of dimension $1$.
Let $a, b \in R$ be nonzerodivisors.
There exists a finite ring extension $R \subset R'$
with $R'/R$ annihilated by a power of $\mathfrak m$
and nonzerodivisors $t, a', b' \in R'$ such that
$a = ta'$ and $b = tb'$ and $R' = a'R' + b'R'$.
\end{lemma}

\begin{proof}
If $a$ or $b$ is a unit, then the lemma is true with $R = R'$.
Thus we may assume $a, b \in \mathfrak m$.
Set $I = (a, b)$. The idea is to blow up $R$ in $I$.
Instead of doing the algebraic argument we work geometrically.
Let $X = \text{Proj}(\bigoplus_{d \geq 0} I^d)$.
By Divisors, Lemma
\ref{divisors-lemma-blowing-up-gives-effective-Cartier-divisor}
the morphism $X \to \Spec(R)$ is an isomorphism over
the punctured spectrum $U = \Spec(R) \setminus \{\mathfrak m\}$.
Thus we may and do view $U$ as an open subscheme of $X$.
The morphism $X \to \Spec(R)$ is projective by
Divisors, Lemma \ref{divisors-lemma-blowing-up-projective}.
Also, every generic point of $X$ lies in $U$, for example
by Divisors, Lemma \ref{divisors-lemma-blow-up-and-irreducible-components}.
It follows from Varieties, Lemma \ref{varieties-lemma-finite-in-codim-1}
that $X \to \Spec(R)$ is finite. Thus $X = \Spec(R')$ is
affine and $R \to R'$ is finite. We have $R_a \cong R'_a$ as $U = D(a)$.
Hence a power of $a$ annihilates the finite $R$-module $R'/R$.
As $\mathfrak m = \sqrt{(a)}$ we see that $R'/R$ is annihilated
by a power of $\mathfrak m$. By
Divisors, Lemma \ref{divisors-lemma-blowing-up-gives-effective-Cartier-divisor}
we see that $IR'$ is a locally principal ideal.
Since $R'$ is semi-local we see that $IR'$ is principal,
see Algebra, Lemma \ref{algebra-lemma-locally-free-semi-local-free},
say $IR' = (t)$. Then we have $a = a't$ and $b = b't$ and everything is
clear.
\end{proof}

\begin{lemma}
\label{lemma-not-infinitely-divisible}
Let $(R, \mathfrak m)$ be a Noetherian local ring of dimension $1$.
Let $a, b \in R$ be nonzerodivisors with $a \in \mathfrak m$.
There exists an integer $n = n(R, a, b)$ such that for a finite ring
extension $R \subset R'$ if $b = a^m c$ for some $c \in R'$, then $m \leq n$.
\end{lemma}

\begin{proof}
Choose a minimal prime $\mathfrak q \subset R$. Observe that
$\dim(R/\mathfrak q) = 1$, in particular $R/\mathfrak q$ is not a field.
We can choose a discrete valuation ring $A$ dominating $R/\mathfrak q$
with the same fraction field, see
Algebra, Lemma \ref{algebra-lemma-dominate-by-dimension-1}. Observe that
$a$ and $b$ map to nonzero elements of $A$ as nonzerodivisors in $R$
are not contained in $\mathfrak q$. Let $v$ be the discrete valuation on $A$.
Then $v(a) > 0$ as $a \in \mathfrak m$.
We claim $n = v(b)/v(a)$ works.

\medskip\noindent
Let $R \subset R'$ be given. Set $A' = A \otimes_R R'$.
Since $\Spec(R') \to \Spec(R)$ is surjective
(Algebra, Lemma \ref{algebra-lemma-integral-overring-surjective})
also $\Spec(A') \to \Spec(A)$ is surjective
(Algebra, Lemma \ref{algebra-lemma-surjective-spec-radical-ideal}).
Pick a prime $\mathfrak q' \subset A'$ lying over $(0) \subset A$.
Then $A \subset A'' = A'/\mathfrak q'$ is a finite extension of rings
(again inducing a surjection on spectra).
Pick a maximal ideal $\mathfrak m'' \subset A''$
lying over the maximal ideal of $A$ and a discrete valuation ring
$A'''$ dominating $A''_{\mathfrak m''}$ (see lemma cited above).
Then $A \to A'''$ is an extension of discrete valuation rings
and we have $b = a^m c$ in $A'''$. Thus $v'''(b) \geq mv'''(a)$.
Since $v''' = ev$ where $e$ is the ramification index
of $A'''/A$, we find that $m \leq n$ as desired.
\end{proof}

\begin{lemma}
\label{lemma-prepare-tame-symbol}
Let $(A, \mathfrak m)$ be a Noetherian local ring of dimension $1$.
Let $r \geq 2$ and let $a_1, \ldots, a_r \in A$ be nonzerodivisors
not all units.
Then there exist
\begin{enumerate}
\item a finite ring extension $A \subset B$ with
$B/A$ annihilated by a power of $\mathfrak m$,
\item for each maximal ideal $\mathfrak m_j \subset B$
a nonzerodivisor $\pi_j \in B_j = B_{\mathfrak m_j}$, and
\item factorizations $a_i = u_{i, j} \pi_j^{e_{i, j}}$ in $B_j$
with $u_{i, j} \in B_j$ units and $e_{i, j} \geq 0$.
\end{enumerate}
\end{lemma}

\begin{proof}
Since at least one $a_i$ is not a unit we find that $\mathfrak m$
is not an associated prime of $A$. Moreover, for any $A \subset B$
as in the statement $\mathfrak m$ is not an associated prime of $B$
and $\mathfrak m_j$ is not an associate prime of $B_j$.
Keeping this in mind will help check the arguments below.

\medskip\noindent
First, we claim that it suffices to prove the lemma for $r = 2$.
We will argue this by induction on $r$; we suggest the reader
skip the proof. Suppose we are given $A \subset B$ and $\pi_j$ in
$B_j = B_{\mathfrak m_j}$ and factorizations
$a_i = u_{i, j} \pi_j^{e_{i, j}}$ for $i = 1, \ldots, r - 1$ in $B_j$
with $u_{i, j} \in B_j$ units and $e_{i, j} \geq 0$.
Then by the case $r = 2$ for $\pi_j$ and $a_r$ in $B_j$
we can find extensions $B_j \subset C_j$ and for every maximal ideal
$\mathfrak m_{j, k}$ of $C_j$ a nonzerodivisor
$\pi_{j, k} \in C_{j, k} = (C_j)_{\mathfrak m_{j, k}}$
and factorizations
$$
\pi_j = v_{j, k} \pi_{j, k}^{f_{j, k}}
\quad\text{and}\quad
a_r = w_{j, k} \pi_{j, k}^{g_{j, k}}
$$
as in the lemma. There exists a unique finite extension $B \subset C$
with $C/B$ annihilated by a power of $\mathfrak m$ such
that $C_j \cong C_{\mathfrak m_j}$ for all $j$, see
Lemma \ref{lemma-glue-at-max}.
The maximal ideals of $C$ correspond $1$-to-$1$
to the maximal ideals $\mathfrak m_{j, k}$ in the localizations
and in these localizations we have
$$
a_i = u_{i, j} \pi_j^{e_{i, j}} =
u_{i, j} v_{j, k}^{e_{i, j}} \pi_{j, k}^{e_{i, j}f_{j, k}}
$$
for $i \leq r - 1$. Since $a_r$ factors correctly too the
proof of the induction step is complete.

\medskip\noindent
Proof of the case $r = 2$. We will use induction on
$$
\ell = \min(\text{length}_A(A/a_1A),\ \text{length}_A(A/a_2A)).
$$
If $\ell = 0$, then either $a_1$ or $a_2$ is a unit and
the lemma holds with $A = B$. Thus we may and do assume $\ell > 0$.

\medskip\noindent
Suppose we have a finite extension of rings $A \subset A'$ such that
$A'/A$ is annihilated by a power of $\mathfrak m$ and such that
$\mathfrak m$ is not an associated prime of $A'$.
Let $\mathfrak m_1, \ldots, \mathfrak m_r \subset A'$
be the maximal ideals and set $A'_i = A'_{\mathfrak m_i}$.
If we can solve the problem for $a_1, a_2$ in each $A'_i$,
then we can apply Lemma \ref{lemma-glue-at-max}
to produce a solution for $a_1, a_2$ in $A$.
Choose $x \in \{a_1, a_2\}$ such that $\ell = \text{length}_A(A/xA)$.
By Lemma \ref{lemma-compare-periodic-lengths} 
and (\ref{equation-multiplicity-coker-ker})
we have $\text{length}_A(A/xA) = \text{length}_A(A'/xA')$.
On the other hand, we have
$$
\text{length}_A(A'/xA') =
\sum [\kappa(\mathfrak m_i) : \kappa(\mathfrak m)]
\text{length}_{A'_i}(A'_i/xA'_i)
$$
by Algebra, Lemma \ref{algebra-lemma-pushdown-module}.
Since $x \in \mathfrak m$ we see that each term on the right hand side
is positive. We conclude that the induction hypothesis applies
to $a_1, a_2$ in each $A'_i$ if $r > 1$ or if $r = 1$ and
$[\kappa(\mathfrak m_1) : \kappa(\mathfrak m)] > 1$.
We conclude that we may assume each $A'$ as above is local with
the same residue field as $A$.

\medskip\noindent
Applying the discussion of the previous paragraph,
we may replace $A$ by the ring constructed in
Lemma \ref{lemma-Noetherian-domain-dim-1-two-elements}
for $a_1, a_2 \in A$. Then since $A$ is local we find,
after possibly switching $a_1$ and $a_2$, that $a_2 \in (a_1)$.
Write $a_2 = a_1^m c$ with $m > 0$ maximal. In fact, by
Lemma \ref{lemma-not-infinitely-divisible}
we may assume $m$ is maximal even after replacing $A$
by any finite extension $A \subset A'$ as in the previous paragraph.
If $c$ is a unit, then we are done. If not, then we replace
$A$ by the ring constructed in
Lemma \ref{lemma-Noetherian-domain-dim-1-two-elements}
for $a_1, c \in A$. Then either (1) $c = a_1 c'$ or
(2) $a_1 = c a'_1$. The first case cannot happen since
it would give $a_2 = a_1^{m + 1} c'$ contradicting the
maximality of $m$. In the second case we get
$a_1 = c a'_1$ and $a_2 = c^{m + 1} (a'_1)^m$.
Then it suffices to prove the lemma for $A$ and $c, a'_1$.
If $a'_1$ is a unit we're done and if not, then
$\text{length}_A(A/cA) < \ell$ because $cA$ is a strictly
bigger ideal than $a_1A$. Thus we win by induction hypothesis.
\end{proof}






